\section{Formal Setting}

To address Q2, a formal framing is required in which the outcome of an enterprise analysis can be evaluated independently of practical preference or organizational expectation. This framing does not introduce definitions or formal machinery; it establishes the conceptual conditions under which the validity of a diagnostic outcome can be assessed.

Any enterprise analysis operates on a model rather than on the enterprise as such. The model specifies a structured space of states, relations, and constraints within which evolution is considered meaningful. Analytical outcomes are therefore not judgments about reality in isolation, but statements about what is admissible within the chosen model. This distinction between model and system follows the classical view of inquiry as an activity mediated by explicitly constructed representations rather than direct access to reality \cite{churchman1971}. Whether continuation is possible can only be assessed relative to this modeled space.

Within this perspective, continuation is not a primitive assumption but a candidate outcome whose validity depends on admissibility. A continuation is admissible if it preserves the internal coherence of the model under its constraints. Conversely, if all candidate continuations violate those constraints, then continuation is not an available outcome of the analysis, regardless of intent or effort. The absence of admissible continuation is not an error of reasoning; it is a property of the modeled system, analogous to constraint-induced limitations in cybernetic systems \cite{ashby1956}.

This distinction separates analytical validity from desirability. Enterprises may strongly prefer forward motion, but preference does not expand the set of admissible states. An analysis that reports the non-existence of admissible continuation does not negate the possibility of action in reality; it states that, under the current model, no internally coherent evolution is representable. The diagnostic content lies precisely in this mismatch between anticipatory expectation and formal representability \cite{rosen1985}.

From this standpoint, STOP or No-Go is neither an exceptional nor a residual category. It occupies the same logical status as any other analytical outcome: it reports the structure of the continuation space as constrained by the model. Treating STOP as failure implicitly assumes that the continuation space is always non-empty if analysis is sufficiently refined. This assumption is not self-evident and therefore requires formal examination.

The framing adopted here shifts attention from decision outcomes to model admissibility. The key question is not whether an analyst could have identified a way forward, but whether a way forward exists within the admissible state space defined by the enterprise model. If no such state exists, then STOP is the only analytically valid result, irrespective of external expectations.

This paper proceeds by analyzing how such situations arise when enterprises are modeled as continua (K\_EA). It investigates how structural constraints, boundaries of admissibility, and internal coherence determine whether the continuation space is non-empty or exhausted. This approach treats the enterprise as a coherent system whose global properties cannot be reduced to isolated local considerations \cite{bertalanffy1968}. The sole aim of this section is to establish the logical framing in which STOP can be evaluated as a diagnostic outcome, without introducing formal definitions or results.
